% Options for packages loaded elsewhere
\PassOptionsToPackage{unicode}{hyperref}
\PassOptionsToPackage{hyphens}{url}
\documentclass[
]{article}
\usepackage{xcolor}
\usepackage{amsmath,amssymb}
\setcounter{secnumdepth}{-\maxdimen} % remove section numbering
\usepackage{iftex}
\ifPDFTeX
  \usepackage[T1]{fontenc}
  \usepackage[utf8]{inputenc}
  \usepackage{textcomp} % provide euro and other symbols
\else % if luatex or xetex
  \usepackage{unicode-math} % this also loads fontspec
  \defaultfontfeatures{Scale=MatchLowercase}
  \defaultfontfeatures[\rmfamily]{Ligatures=TeX,Scale=1}
\fi
\usepackage{lmodern}
\ifPDFTeX\else
  % xetex/luatex font selection
\fi
% Use upquote if available, for straight quotes in verbatim environments
\IfFileExists{upquote.sty}{\usepackage{upquote}}{}
\IfFileExists{microtype.sty}{% use microtype if available
  \usepackage[]{microtype}
  \UseMicrotypeSet[protrusion]{basicmath} % disable protrusion for tt fonts
}{}
\makeatletter
\@ifundefined{KOMAClassName}{% if non-KOMA class
  \IfFileExists{parskip.sty}{%
    \usepackage{parskip}
  }{% else
    \setlength{\parindent}{0pt}
    \setlength{\parskip}{6pt plus 2pt minus 1pt}}
}{% if KOMA class
  \KOMAoptions{parskip=half}}
\makeatother
\usepackage{longtable,booktabs,array}
\newcounter{none} % for unnumbered tables
\usepackage{calc} % for calculating minipage widths
% Correct order of tables after \paragraph or \subparagraph
\usepackage{etoolbox}
\makeatletter
\patchcmd\longtable{\par}{\if@noskipsec\mbox{}\fi\par}{}{}
\makeatother
% Allow footnotes in longtable head/foot
\IfFileExists{footnotehyper.sty}{\usepackage{footnotehyper}}{\usepackage{footnote}}
\makesavenoteenv{longtable}
\usepackage{graphicx}
\makeatletter
\newsavebox\pandoc@box
\newcommand*\pandocbounded[1]{% scales image to fit in text height/width
  \sbox\pandoc@box{#1}%
  \Gscale@div\@tempa{\textheight}{\dimexpr\ht\pandoc@box+\dp\pandoc@box\relax}%
  \Gscale@div\@tempb{\linewidth}{\wd\pandoc@box}%
  \ifdim\@tempb\p@<\@tempa\p@\let\@tempa\@tempb\fi% select the smaller of both
  \ifdim\@tempa\p@<\p@\scalebox{\@tempa}{\usebox\pandoc@box}%
  \else\usebox{\pandoc@box}%
  \fi%
}
% Set default figure placement to htbp
\def\fps@figure{htbp}
\makeatother
\ifLuaTeX
  \usepackage{luacolor}
  \usepackage[soul]{lua-ul}
\else
  \usepackage{soul}
\fi
\setlength{\emergencystretch}{3em} % prevent overfull lines
\providecommand{\tightlist}{%
  \setlength{\itemsep}{0pt}\setlength{\parskip}{0pt}}
\usepackage{bookmark}
\IfFileExists{xurl.sty}{\usepackage{xurl}}{} % add URL line breaks if available
\urlstyle{same}
\hypersetup{
  hidelinks,
  pdfcreator={LaTeX via pandoc}}

\author{}
\date{}

\begin{document}

\begin{quote}
KYAMBOGO
\includegraphics[width=1.24653in,height=0.97902in]{./media/image1.png}UNIVERSITY

FACULTY OF ENGINEERING

BACHELOR OF ENGINEERING IN BIOMEDICAL AND MECHATRONICS ENGINEERING

YEAR ONE: SEMESTER TWO

TEMB 1202: COMPUTING II

TITLE: LIBRARY BOOK MANAGEMENT SYSTEM USING C PROGRAMMING

DATE: 26\textsuperscript{TH} MAY, 2026

GROUP ONE EVENING

\textbf{\ul{MEMBERS}}
\end{quote}

{\def\LTcaptype{none} % do not increment counter
\begin{longtable}[]{@{}
  >{\raggedright\arraybackslash}p{(\linewidth - 2\tabcolsep) * \real{0.3283}}
  >{\centering\arraybackslash}p{(\linewidth - 2\tabcolsep) * \real{0.4168}}@{}}
\toprule\noalign{}
\begin{minipage}[b]{\linewidth}\raggedright
\begin{quote}
\textbf{NAMES}
\end{quote}
\end{minipage} & \begin{minipage}[b]{\linewidth}\centering
\begin{quote}
\textbf{REGISTRATION NUMBER}
\end{quote}
\end{minipage} \\
\midrule\noalign{}
\endhead
\bottomrule\noalign{}
\endlastfoot
\begin{minipage}[t]{\linewidth}\raggedright
\begin{quote}
OREM ABRAHAM
\end{quote}
\end{minipage} & \begin{minipage}[t]{\linewidth}\centering
\begin{quote}
25/U/BIE/01415/PE
\end{quote}
\end{minipage} \\
\begin{minipage}[t]{\linewidth}\raggedright
\begin{quote}
JUUKO MORRIS
\end{quote}
\end{minipage} & \begin{minipage}[t]{\linewidth}\centering
\begin{quote}
25/U/BIE/01375/PE
\end{quote}
\end{minipage} \\
\begin{minipage}[t]{\linewidth}\raggedright
\begin{quote}
OKOT CLEMENT
\end{quote}
\end{minipage} & \begin{minipage}[t]{\linewidth}\centering
\begin{quote}
25/U/BIE/05306/PE
\end{quote}
\end{minipage} \\
\begin{minipage}[t]{\linewidth}\raggedright
\begin{quote}
SANYU DORCUS
\end{quote}
\end{minipage} & \begin{minipage}[t]{\linewidth}\centering
\begin{quote}
25/U/BIE/01420/PE
\end{quote}
\end{minipage} \\
\begin{minipage}[t]{\linewidth}\raggedright
\begin{quote}
EDYABI BRIAN
\end{quote}
\end{minipage} & \begin{minipage}[t]{\linewidth}\centering
\begin{quote}
25/U/BIE/05291/PE
\end{quote}
\end{minipage} \\
\begin{minipage}[t]{\linewidth}\raggedright
\begin{quote}
BUYONDO EMMANUEL
\end{quote}
\end{minipage} & \begin{minipage}[t]{\linewidth}\centering
\begin{quote}
25/U/BIO/04517/PE
\end{quote}
\end{minipage} \\
\end{longtable}
}

\begin{enumerate}
\def\labelenumi{\arabic{enumi}.}
\item
  \textbf{Introduction}
\end{enumerate}

This report documents the design and implementation of a Library Book
Management System, developed as part of the TEMB 1202 Computing II group
assignment. Written entirely in the C programming language, the system
provides a full-featured, menu-driven console application for managing a
library\textquotesingle s book inventory, borrowing workflow, fine
collection, and reporting. Libraries are vital knowledge hubs, and
effective management of their resources directly impacts user
satisfaction and institutional efficiency. Manual record-keeping leads
to lost books, uncollected fines, and inaccurate availability
information. This system automates these processes using structured C
programming techniques structs, arrays, functions, and conditional logic
to enforce library policies consistently.

Scope: The system manages up to 100 book records and 50 borrower
records, enforces five business rules, tracks borrow counts to flag
high-demand books, calculates late-return fines, blocks borrowing for
members with unpaid fines, and generates a comprehensive library summary
report identifying the most-borrowed book.

\begin{enumerate}
\def\labelenumi{\arabic{enumi}.}
\setcounter{enumi}{1}
\item
  \textbf{The problem being solved}
\end{enumerate}

University and community libraries need a reliable, automated system to
track book availability, manage borrower records, enforce borrowing
policies, and calculate fines for late returns. The assignment requires
a C program that addresses all these needs while applying a set of
clearly defined business rules.

\begin{enumerate}
\def\labelenumi{\alph{enumi})}
\item
  Objectives
\end{enumerate}

\begin{itemize}
\item
  Allow authorised librarians to add new book records with full
  bibliographic details.
\item
  Enable searching for books by ID, title, or author.
\item
  Support updating book details and deleting records.
\item
  Process borrow transactions, enforcing availability and fine-status
  rules.
\item
  Process return transactions, calculating and recording fines for late
  returns.
\item
  Prevent borrowing when a member has unpaid outstanding fines.
\item
  Flag books that have been borrowed many times as high-demand.
\item
  Generate a summary report showing borrow counts, fines collected, and
  most-borrowed book.
\item
  Track total borrowed books and available copies in real time.
\item
  Require authentication before system access.
\end{itemize}

\begin{enumerate}
\def\labelenumi{\arabic{enumi}.}
\setcounter{enumi}{2}
\item
  \textbf{System features}
\end{enumerate}

The system provides eight operations accessible through a main menu,
organised into six functional modules. Each module is implemented as one
or more dedicated C functions.

{\def\LTcaptype{none} % do not increment counter
\begin{longtable}[]{@{}
  >{\raggedright\arraybackslash}p{(\linewidth - 4\tabcolsep) * \real{0.2556}}
  >{\raggedright\arraybackslash}p{(\linewidth - 4\tabcolsep) * \real{0.2605}}
  >{\raggedright\arraybackslash}p{(\linewidth - 4\tabcolsep) * \real{0.2555}}@{}}
\toprule\noalign{}
\begin{minipage}[b]{\linewidth}\raggedright
Feature
\end{minipage} & \begin{minipage}[b]{\linewidth}\raggedright
Description
\end{minipage} & \begin{minipage}[b]{\linewidth}\raggedright
Rules applied
\end{minipage} \\
\midrule\noalign{}
\endhead
\bottomrule\noalign{}
\endlastfoot
Add Book & Register a new book: title, author, ISBN, genre, copies, year
& Duplicate ISBN check \\
Search Book & Find by ID, title keyword, or author name. &
Case-insensitive match \\
Update Book & Modify any field of an existing book record. & ID must
exist \\
Delete Book & Remove a book record with confirmation. & Cannot delete if
copies borrowed \\
Borrow Book & Lend a copy to a member; records due date. & BR-1, BR-2,
BR-4 enforced \\
Return Book & Accept a returned copy; compute and record fine if late. &
BR-3 enforced \\
Display Books & List all books with availability and demand status &
BR-5 high-demand flag \\
Library Report & Aggregate borrow counts, fines, and identify
most-borrowed book. & Full summary \\
\end{longtable}
}

\begin{enumerate}
\def\labelenumi{\alph{enumi})}
\item
  Book Registration \& Management
\end{enumerate}

\begin{quote}
When adding a book the user supplies: a unique numeric ID, ISBN, title,
author, genre, publication year, and total number of copies. The system
validates that the ID and ISBN are not duplicates and that all numeric
values are positive before saving the record. The available-copies count
is initialised to the total copies value.
\end{quote}

\begin{enumerate}
\def\labelenumi{\alph{enumi})}
\setcounter{enumi}{1}
\item
  Search \& Display
\end{enumerate}

\begin{quote}
Three search modes are supported: by book ID (exact match), by title
keyword (substring search), and by author name (substring search). The
display-all view presents every book with its availability count and
automatically annotates high-demand books with a visible flag.
\end{quote}

\begin{enumerate}
\def\labelenumi{\alph{enumi})}
\setcounter{enumi}{2}
\item
  Borrow Books
\end{enumerate}

\begin{quote}
A borrow transaction requires a member ID and book ID. The system checks
copy availability (BR-1/BR-2), verifies the member has no unpaid fines
(BR-4), records the borrow with a due date 14 days from today,
decrements available copies, increments the borrow count, and flags the
book as high-demand if the count exceeds the threshold (BR-5).
\end{quote}

\begin{enumerate}
\def\labelenumi{\alph{enumi})}
\setcounter{enumi}{3}
\item
  Return Books \& Fine Calculation
\end{enumerate}

\begin{quote}
On return the system computes the number of days overdue. If the book is
late (BR-3), a fine of 500 UGX per overdue day is calculated and added
to the member\textquotesingle s outstanding balance. The member cannot
borrow another book until the fine is cleared (BR-4). Available copies
are incremented on successful return.
\end{quote}

\begin{enumerate}
\def\labelenumi{\alph{enumi})}
\setcounter{enumi}{4}
\item
  Library Summary Report
\end{enumerate}

\begin{quote}
The report lists all books with their total borrow count and remaining
copies, prints cumulative statistics (total borrows, total fines
collected), and identifies the most-borrowed book, giving management an
at-a-glance view of library performance.
\end{quote}

\begin{enumerate}
\def\labelenumi{\arabic{enumi}.}
\setcounter{enumi}{3}
\item
  \textbf{Program design and structure}
\end{enumerate}

\begin{enumerate}
\def\labelenumi{\alph{enumi})}
\item
  Data Structures
\end{enumerate}

\begin{quote}
Two structs model the library domain: Book (catalogue record) and Member
(borrower record). Both are stored in fixed-size arrays.
\end{quote}

\begin{enumerate}
\def\labelenumi{\alph{enumi})}
\setcounter{enumi}{1}
\item
  System Architecture
\end{enumerate}

\begin{quote}
The program follows a three-layer architecture: UI (menus, formatted
output), Business Logic (rules, calculations), and Data (struct arrays,
(Create, Read, Update, Delete) CRUD helpers).
\end{quote}

{\def\LTcaptype{none} % do not increment counter
\begin{longtable}[]{@{}
  >{\raggedright\arraybackslash}p{(\linewidth - 4\tabcolsep) * \real{0.2349}}
  >{\raggedright\arraybackslash}p{(\linewidth - 4\tabcolsep) * \real{0.2348}}
  >{\raggedright\arraybackslash}p{(\linewidth - 4\tabcolsep) * \real{0.2348}}@{}}
\toprule\noalign{}
\begin{minipage}[b]{\linewidth}\raggedright
Layer
\end{minipage} & \begin{minipage}[b]{\linewidth}\raggedright
Responsibility
\end{minipage} & \begin{minipage}[b]{\linewidth}\raggedright
Key Functions
\end{minipage} \\
\midrule\noalign{}
\endhead
\bottomrule\noalign{}
\endlastfoot
UI / Menu & Display menus, format tables, collect validated input &
main(), displayMenu(), printBookTable() \\
Business logic & Enforce all 5 rules, compute fines, set demand flags &
borrowBook(), returnBook(),

calculateFine() \\
Data management & CRUD on Book and Member arrays & addBook(),
searchBook(), updateBook(), deleteBook( \\
Reporting & Aggregate statistics, find most-borrowed book &
generateReport(), findMostBorrowed() \\
Authentication & Verify librarian credentials at startup &
authenticate() \\
\end{longtable}
}

\begin{enumerate}
\def\labelenumi{\alph{enumi})}
\setcounter{enumi}{2}
\item
  \textbf{Module Breakdown}
\end{enumerate}

\begin{itemize}
\item
  Authentication Module: Simple three-attempt username/password gate.
  Hardcoded for demo; easily extended to file-based storage.
\item
  Book Management Module: addBook(), searchById(), searchByTitle(),
  searchByAuthor(), updateBook(), deleteProduce(). All operations
  validate input and provide clear error messages.
\item
  Member Management Module: registerMember() collects member details;
  findMemberById() is used by borrow/return operations.
\item
  Borrow Module: borrowBook() enforces BR-1, BR-2, BR-4, BR-5, records
  the loan, and updates counts and dates.
\item
  Return Module: returnBook() computes days overdue, invokes
  calculateFine(), updates the member\textquotesingle s unpaid balance,
  and restores copy count.
\item
  Reporting Module: generateReport() iterates both arrays, accumulates
  totals, and calls findMostBorrowed() for the top result.
\item
  Utility Module: clearBuffer(), computeDaysDiff(), dateToInt() ---
  helpers for safe input and date arithmetic.
\end{itemize}

\begin{enumerate}
\def\labelenumi{\arabic{enumi}.}
\setcounter{enumi}{4}
\item
  \textbf{Important function used.}
\end{enumerate}

\begin{itemize}
\item
  Authentication: The authenticate() function is called once at startup,
  granting three attempts before the program exits. Default credentials
  are admin / lib2026.
\end{itemize}

\begin{quote}
\includegraphics[width=6.26806in,height=4.08403in]{./media/image2.png}
\end{quote}

\begin{itemize}
\item
  Book Management --- addBook(): Collects bibliographic fields, checks
  for duplicate ID/ISBN, then appends the new Book struct.
  availableCopies is set equal to totalCopies; borrowCount and
  isHighDemand are initialised to zero.
\end{itemize}

\begin{quote}
\includegraphics[width=5.80972in,height=3.89868in]{./media/image3.png}
\end{quote}

\begin{itemize}
\item
  Borrow \& Return Logic
\end{itemize}

\begin{quote}
\includegraphics[width=6.26806in,height=5.90903in]{./media/image4.png}
\end{quote}

\begin{itemize}
\item
  Fine Calculation --- returnBook()
\end{itemize}

\begin{quote}
\includegraphics[width=5.92431in,height=3.64936in]{./media/image5.png}
\end{quote}

\begin{itemize}
\item
  Reporting --- generateReport()
\end{itemize}

\begin{quote}
\includegraphics[width=5.96597in,height=5.86022in]{./media/image6.png}
\end{quote}

\begin{enumerate}
\def\labelenumi{\arabic{enumi}.}
\setcounter{enumi}{5}
\item
  \textbf{Design decisions and justifications}
\end{enumerate}

\begin{itemize}
\item
  \textbf{Two-Struct Design}: Separating book data (Book) from borrower
  data (Member) reflects real library systems and avoids redundant data
  storage. Each struct is cohesive and independently extensible.
\item
  \textbf{One-Book-Per-Member Limit}: A member may only borrow one book
  at a time in the current implementation. This simplifies the Member
  struct and prevents complex multi-loan tracking while still meeting
  all stated assignment requirements. Multiple borrows could be added
  using a linked list in a future version.
\item
  \textbf{Parameterised Constants via \#define}: Fine rate, loan period,
  and high-demand threshold are all \#define macros. This means
  institutional policy changes (e.g. changing the loan period from 14 to
  21 days) require editing a single line rather than hunting through the
  code.
\item
  \textbf{Fine Blocking at Borrow Time (BR-4):} The fine check occurs at
  the beginning of borrowBook(), before any other processing. This
  ensures the rule is enforced consistently and the system provides a
  clear, actionable message telling the member how much they owe.
\item
  \textbf{High-Demand Flag (BR-5):} The isHighDemand field is set lazily
  --- updated each time borrowCount crosses the threshold. This avoids
  scanning the entire catalogue on every display call and keeps the flag
  persistently accurate without additional computation.
\item
  \textbf{Pay Fine Feature:} A dedicated payFine() function was added
  because the borrow-block rule (BR-4) is otherwise irresolvable within
  the system. This makes the application self-contained and realistic.
\item
  \textbf{Simplified Date Arithmetic:} Full calendar date arithmetic
  (accounting for leap years and month lengths) was simplified to a
  linear approximation suitable for the demonstration scope. A
  production system would use mktime() and difftime() from.
\end{itemize}

\begin{enumerate}
\def\labelenumi{\arabic{enumi}.}
\setcounter{enumi}{6}
\item
  \textbf{Challenges encountered}
\end{enumerate}

\begin{itemize}
\item
  \textbf{Managing Two Interdependent Structs:} The borrow and return
  operations must consistently update both the Book and Member structs
  atomically. A bug that updated one but not the other would corrupt
  availability counts. This was handled by always finding both records
  before making any modifications.
\item
  \textbf{Input Buffer Flushing:} Mixing scanf() with fgets() required
  careful buffer management. A clearBuffer() helper consuming the
  newline after every numeric scanf call resolved all input-skipping
  issues.
\item
  \textbf{Preventing Invalid Delete Operations:} Deleting a book while
  copies are on loan would create dangling borrowedBookId references in
  Member records. The delete function checks availableCopies ==
  totalCopies before allowing deletion, preventing this data integrity
  issue.
\item
  \textbf{Date Storage and Comparison:} Storing dates as strings
  (DD/MM/YYYY) is human-readable but makes arithmetic harder. The
  computeOverdueDays() helper converts the string to an approximate
  integer for comparison. A real system would use time, t for reliable
  date handling.
\item
  \textbf{Partial Fine Payments:} The assignment only states that unpaid
  fines prevent borrowing. Adding partial payment support required
  careful thinking about when to clear the block (only when balance
  reaches zero) versus allowing partial payments to reduce the balance.
\end{itemize}

\begin{enumerate}
\def\labelenumi{\arabic{enumi}.}
\setcounter{enumi}{7}
\item
  \textbf{Conclusion}
\end{enumerate}

The Library Book Management System successfully fulfils every
requirement of the TEMB 1202 Computing II Group A Evening assignment.
The system delivers a complete, menu-driven C application that manages
books and members through their full lifecycle from catalogue
registration to borrowing, return, fine collection, and summary
reporting.

All five business rules are correctly enforced: unavailable books cannot
be borrowed, an explicit unavailability message is shown when copies are
exhausted, late returns trigger a per-day fine, members with unpaid
fines are blocked from borrowing, and frequently borrowed

books are automatically flagged as high demand.

The program demonstrates proficient use of C programming concepts
including structures, array, pointers, string manipulation, functions,
and conditional logic --- all topics covered in the Computing II course
unit.

Potential future enhancements include file-based persistence for data
across sessions, support for multiple simultaneous borrows per member,
full calendar-accurate date arithmetic using, and a role-based access
control system distinguishing librarian administrators from general
staff.

\end{document}
